\documentclass[12pt]{article}

\usepackage{amssymb}
\usepackage{amsmath}
\usepackage{bm}
\usepackage{color}
\usepackage{float}
\usepackage{graphicx}
\usepackage{mathtools}
\usepackage[none]{hyphenat}
\usepackage{indentfirst}
\usepackage[labelfont=bf]{caption}
\usepackage[left=2.5cm,right=2.5cm, top=2.5cm,bottom=2.5cm]{geometry}
\usepackage{tcolorbox}
\usepackage{lineno}
\usepackage{setspace}
\DeclarePairedDelimiter\ceil{\lceil}{\rceil}

\usepackage[authoryear]{natbib}
\usepackage{hyperref}
\usepackage{yhmath}
\usepackage{accents}

\usepackage{dsfont}
\usepackage{mathtools}
\usepackage{soul}

\overfullrule 5pt

\doublespacing

\usepackage{lipsum}

%\usepackage{adjustbox}
\usepackage{lscape}
\usepackage{afterpage}

\makeatletter
\renewcommand{\maketitle}{\bgroup\setlength{\parindent}{0pt}
\begin{flushleft}
\textbf{\@title}

  \@author
\end{flushleft}\egroup
}
\makeatother

\begin{document}

\title{Seafood Fraud Social Food Web Model}

\maketitle

The below deterministic discrete-time Consumer-Resource model is inspired by those of \citet{fryxell2010resource}, \citet{mccann2016food}, \citet{fryxell2017supply}, \citet{bieg2017dynamical}, and \citet{bieg2023fisheries}, incorporating elements of fisheries bioeconomic theory, community ecology theory, and human behavioural ecology theory. 

The Resource (Seafood biomass; $S$) follows a Ricker logistic surplus production dynamics model of net recruitment subject to density dependence (through maximum per capita exponential growth rate ($r_{\mathrm{max}, S}$) and linear harvest (\textit{i.e.}, catch) ($H_S$)). Harvest comes from two types of Consumers: Buyers ($B$), who seek to purchase a proportion ($f$) of legally-caught and correctly-labelled seafood at market price, as well as Fraudsters ($F$), who seek to illicitly misrepresent the remaining fraction ($1 - f$) of illegally-caught seafood through adulteration, mislabelling, substitution, or other forms of fraud. In this sense, Seafood can be viewed as producers, Fraudsters as predators, and Buyers as prey, interacting with one another in a social meta food web (\textit{i.e.}, the seafood supply chain) \citep{bieg2022replicating} linked by smaller food webs (\textit{i.e.}, fisheries) that exchange information \citep{bieg2018linking}.  Buyers and Fraudsters compete with one another for the Resource and their levels are regulated through their own behaviours and profits. Harvest of seafood depends on fishing effort/intensity ($E$) expended to obtain seafood (\textit{e.g.}, the number of boats, the number of fishers, the number of meters of gillnet, the number of hooks set, \textit{etc.}) scaled by catchability ($q$) of the seafood (the ability of Buyers and Fraudsters to obtain the resource for every one unit of expended effort). Since Buyers and Fraudsters make up the entire human population with the system, their levels are conserved; thus, their levels at time $t$ (as proportions of their carrying capacities, which are set to one) must sum to one (\textit{i.e.}, $B_t$ + $F_t$ = 1). Therefore,

\begin{align}
B^{'}_{t} &= \frac{B_{t}}{B_{t} + F_{t}} \\
F^{'}_{t} &= \frac{F_{t}}{B_{t} + F_{t}} = 1 - B^{'}_{t}.
\end{align}

Harvest of Seafood by Buyers and Fraudsters, respectively, is given by 

\begin{align}
H_{B, t} &= q_{B}E_{B, t}S_t \\
H_{F, t} &= q_{F}E_{F, t}S_t
\end{align}

\noindent where harvesting Effort by Buyers and Fraudsters, respectively, takes the form

\begin{align}
E_{B, t} &= f B_t \\
E_{F, t} &= (1 - f) F_t
\end{align}

\noindent so that,

\begin{align}
H_{B, t} &= q_{B}f B_tS_t \\
H_{F, t} &= q_{F}(1 - f) F_tS_t
\end{align}

where the total Harvest of Seafood is

\begin{align}
H_{\mathrm{S}, t} &= H_{B, t} + H_{F, t} \\
 &= q_{B}f B_tS_t + q_{F}(1 - f) F_tS_t.
\end{align}

Therefore, the Seafood equation is 

\begin{equation}
S_{t+1} = S_te^{r_{\mathrm{max}, S}(1 - S_t)} - q_{B}f B_tS_t - q_{F}(1 - f)F_tS_t.
\end{equation}

Note, the Catch-per-unit-Effort (CPUE) by Buyers and Fraudsters, respectively, is defined as

\begin{align}
\mathrm{CPUE}_{B, t} &= q_{B}S_t \\
\textrm{CPUE}_{F, t} &= q_{F}S_t
\end{align}

\noindent and the fishing ``mortality" (\textit{i.e.}, harvest/catch divided by biomass) due to purcurment of seafood by Buyers and Fraudsters, respectively, is

\begin{align}
m_{B, t} &= q_{B} B_t \\
m_{F, t} &= q_{F} F_t.
\end{align}

Equation (11) can thus be expressed as

\begin{align}
S_{t+1} &= S_te^{r_{\mathrm{max}, S}(1 - S_t)} - fB_t\mathrm{CPUE}_{B, t} - (1 - f)F_t\textrm{CPUE}_{F, t} \\
&= S_te^{r_{\mathrm{max}, S}(1 - S_t)} - fm_{B, t}S_t - (1 - f)m_{F, t}S_t.
\end{align}

Fraudsters become more irrational as the difference between market seafood price and wholesale seafood price increases ($\Delta P = P_{\textrm{Market}} - P_{\textrm{Wholesale}}$). Typically, market seafood prices are higher than wholesale prices due to the added processing of seafood required prior to the point of sale. When $\Delta P < 0$, Fraudsters are negatively impacted, decreasing their levels. In the case of $\Delta P > 0$, Fraudsters benefit, increasing their levels. If $\Delta P = 0$, Fraudster levels remain unchanged.

Buyers and Fraudsters seek to optimize their overall behaviours and profits in obtaining Seafood. Profit is the difference between revenue of the seafood (landed price times harvest) and total cost of the seafood (cost per unit effort times effort). That is,

\begin{align}
\Pi_{B, t} &= \pi_B q_{B}f B_tS_t - c_B f B_t \\
&= \pi_B f B_t \mathrm{CPUE}_{B, t} - c_B f B_t \\
&= \pi_B fm_{B, t}S_t  - c_B  f B_t \\
\Pi_{F, t} &= \pi_F q_{F}(1 - f) F_tS_t - c_F (1 -f) F_t \\
&= \pi_F (1 - f) F_t \mathrm{CPUE}_{B, t} - c_F  (1 - f) F_t \\
&= \pi_F (1 - f) m_{F, t}S_t  - c_F (1 - f) F_t .
\end{align}

The Buyer equation is thus given by

\begin{align}
B_{t+1} &= B_te^{\alpha B_t S_t - \beta B_t F_t + \pi_B q_{B}f B_tS_t - c_B f B_t} \\
&= B_te^{\alpha B_t S_t - \beta B_t F_t + \pi_B f B_t \mathrm{CPUE}_{B, t} - c_B f B_t} \\ 
&= B_te^{\alpha B_t S_t - \beta B_t F_t \pi_B fm_{B, t}S_t  - c_B  f B_t}.
\end{align}

The Fraudster equation takes the form

\begin{align}
F_{t+1} &= F_te^{\big(\gamma F_t S_t - \delta F_t B_t\big)\epsilon\Delta P + \pi_F q_{F}(1 - f) F_tS_t - c_F (1 - f) F_t}  \\
&= F_te^{\big(\gamma F_t S_t - \delta F_t B_t\big)\epsilon\Delta P + \pi_F (1 - f) F_t \mathrm{CPUE}_{B, t} - c_F (1 - f) F_t} \\
&= F_te^{\big(\gamma F_t S_t - \delta F_t B_t\big)\epsilon\Delta P + \pi_F (1 - f) m_{F, t} S_t  - c_F (1 - f) F_t}.
\end{align}

The parameters $\alpha$, $\beta$, $\gamma$, and $\delta$ represent the degree of awareness of Buyers to changing levels of seafood, the degree of awareness of Buyers to evolving Fraudster behaviour, the degree of awareness of Fraudsters to shifting seafood levels, and the degree of awareness of Fraudsters to altering Buyer behaviour, respectively. The parameter $\epsilon$ reflects the degree of awareness of Fraudsters to price volatility.

The quantities $\alpha B_t S_t - \beta B_t F_t$ and ($\gamma F_t S_t - \delta F_t B_t\big)\epsilon\Delta P$ in Equations (24)-(26) and Equations (27)-(29) can be thought of as behaviour terms for Buyers and Fraudsters, \\ respectively. The quantities $\pi_B q_{B}f B_tS_t - c_BfB_t$ and $\pi_F q_{F}(1 - f) F_tB_t - c_F (1 - f) F_t$ within the same equations are profitability terms, where $\pi_B$ and $\pi_F$ are prices for seafood paid by Buyers and Fraudsters, respectively, and $c_B$ and $c_F$ are costs per unit effort expended by Buyers and Fraudsters to obtain seafood, respectively.

The final model contains three variables indexed by time ($S_t$, $B_t$, and $F_t$), 14 fixed parameters ($r_{\mathrm{max}, S}$, $f$, $q_B$, $q_F$, $\pi_B$, $\pi_F$, $c_B$, $c_F$, $\alpha$, $\beta$, $\gamma$, $\delta$, $\epsilon$, and $\Delta P$), and four derived quantities ($\mathrm{CPUE}_{B, t}$, $\mathrm{CPUE}_{F, t}$, $m_{B, t}$, and $m_{F, t}$) . Due to the nonlinear nature of the model, all parameters except $r_{\mathrm{max}, S}$ ($>$ 0), and $\Delta P$ are constrained to [0, 1] to avoid very large, very small, or infinite values of the variables that could arise due to the exponential terms within the model. 


\newpage

\section*{Equilibria and Stability}

The system has eight equilibria (with many special cases):

\begin{align}
(S^*, B^*, F^*) &= (0, 1, 0) \\
(S^*, B^*, F^*) &= (0, 0, 1) \\
(S^*, B^*, F^*) &= (1, 1, 0) \\
(S^*, B^*, F^*) &= (1, 0, 1) \\
(S^*, B^*, F^*) &= \Big(1 - \frac{1}{r_{\mathrm{max}, S}}\mathrm{ln}(1 + q_Bf), 1, 0\Big) \\
(S^*, B^*, F^*) &= \Big(1 - \frac{1}{r_{\mathrm{max}, S}}\mathrm{ln}(1 + q_F(1 - f)), 0, 1\Big)
\\
(S^*, B^*, F^*) &= \Big(\frac{\beta + c_Bf}{\alpha + \pi_B q_Bf}, 0, 1\Big)
\\
(S^*, B^*, F^*) &= \Big(\frac{\delta \epsilon \Delta P + c_F(1 - f)}{\gamma \epsilon \Delta P + \pi_F q_F(1 -f)}, 1, 0\Big).
\end{align}

\noindent When $q_B$ = 0, $f$ = 0, or both, with $r_{\mathrm{max}, S} > 0$ in all cases, Equation (34) becomes

\begin{equation}
(S^*, B^*, F^*) = (1, 1, 0).
\end{equation}

\noindent When both $q_B$ = 1 and $f$ = 1, with $r_{\mathrm{max}, S} > 0$ in all cases, Equation (34) becomes

\begin{equation}
(S^*, B^*, F^*) = (1 - \mathrm{ln}(2), 1, 0) \approx (0.307, 1, 0).
\end{equation}

\noindent When $q_F$ = 1 and $f$ = 0, with $r_{\mathrm{max}, S}$ = 1 in all cases, Equation (35) becomes

\begin{equation}
(S^*, B^*, F^*) = (1 - \mathrm{ln}(2), 0, 1) \approx (0.307, 0, 1).
\end{equation}

\noindent Similarly, when $q_F$ = 0 and $f$ = 1, with $r_{\mathrm{max}, S}$ = 1 in all cases, Equation (35) becomes

\begin{equation}
(S^*, B^*, F^*) = (1 - \mathrm{ln}(2), 0, 1) \approx (0.307, 0, 1).
\end{equation}

\noindent When $f$ = 0, Equation (36) becomes

\begin{equation}
(S^*, B^*, F^*) = \Big(\frac{\beta}{\alpha}, 0, 1\Big),
\end{equation}

\noindent with the added condition that $\alpha \geq \beta$ to ensure $S^* \in [0, 1]$.

\noindent When both $\alpha$ = 0 and $\beta$ = 0, Equation (36) becomes

\begin{equation}
(S^*, B^*, F^*) = \Big(\frac{c_B}{\pi_Bq_B}, 0, 1\Big),
\end{equation}

\noindent for any value of $0 < f \leq 1$, with the additional contraint that $c_B \leq \pi_B q_B$ to ensure $S^* \in [0, 1]$.

\noindent When $f$ = 1, Equation (37) becomes

\begin{equation}
(S^*, B^*, F^*) = \Big(\frac{\delta}{\gamma}, 1, 0\Big),
\end{equation}

\noindent for any value of $\epsilon \neq 0$ and $\Delta P \neq 0$, with the additional constraint that $\delta \leq \gamma$ to ensure $S^* \in [0, 1]$.

\noindent When one or more of $\gamma$, $\delta$, $\epsilon$, and $\Delta P$ is zero, Equation (37) reduces to 

\begin{equation}
(S^*, B^*, F^*) = \Big(\frac{c_F}{\pi_Fq_F}, 1, 0\Big),
\end{equation}

\noindent for any value of $0 \leq f < 1$, with the additional constraint that $c_F \leq \pi_F q_F$ to ensure $S^* \in [0, 1]$.


From Equations (36) and (37), a sharper bound on $f$ is

\begin{equation}
f \in \Big[\frac{\beta - \alpha}{\pi_B q_B - c_B}, \frac{(\gamma - \delta)\epsilon}{\pi_F q_F - c_F}\Delta P + 1 \Big]
\end{equation}

\noindent which reduces to $f \in [0, 1]$ (which is the case for Equations (32) and (33)) when $\alpha = \beta$ and either $\gamma = \delta$, $\epsilon = 0$, or $\Delta P = 0.$

Are there any biologically meaningful combinations of variables/parameters leading to oscillations, bifurcations, limit cycles, or chaos?

\section*{Model Extensions}

\begin{itemize}

\item Add interaction term between Seafood, Buyers and Fraudsters to Buyer and Fraudster equations -- $\eta B_tF_tS_t$ (not sure of the sign, positive or negative?)

\vspace{1mm}

\item Wild \textit{vs.} farmed fish

\vspace{1mm}

\item Ethical harvesters (\textit{e.g.}, Indigenous communities)

\vspace{1mm}

\item Legislators/policymakers/fisheries managers

\vspace{1mm}

\item Agent-based model - Seafood, Buyer, and Fraudster levels are updated according to simple deterministic/stochastic rules

\vspace{1mm}

\item Game theory - Buyers and Fraudsters play a zero-sum game as they compete for Seafood, whereby one player's gain is the other's loss

\end{itemize}

\newpage

\bibliographystyle{humannat}
\bibliography{References}

\end{document}